\chapter{Pelvic Examination}

\section*{Overview}
A pelvic examination assesses the external and internal reproductive organs, including the vulva, vagina, cervix, uterus, and adnexa. It may be part of routine care or an evaluation of symptoms.

\section*{Why It Is Done}
It can help evaluate pelvic pain, abnormal bleeding, discharge, lesions, urinary or sexual symptoms, pregnancy-related concerns, and changes identified during screening. Samples for cervical screening or infection testing may be collected when appropriate.

\section*{How It Is Performed}
After discussion and consent, the patient lies on an examination table. The clinician inspects the external area, may use a speculum to view the vaginal walls and cervix, and may perform a gentle bimanual examination to assess the uterus and adnexa. The patient may ask to stop at any time.

\section*{Preparation and Risks}
No special preparation is usually needed. Emptying the bladder may improve comfort. Pressure or brief discomfort can occur, but severe pain should be reported. Minor spotting may occur after sampling.

\section*{Outcome and Follow-up}
Findings may be normal or may suggest infection, inflammation, structural changes, or a need for additional testing. A pelvic examination does not replace cervical cancer screening or other tests when those are indicated.

\section*{References}
\begin{enumerate}
  \item MedlinePlus. Pelvic examination. U.S. National Library of Medicine; 2025.
  \item Current Medical Diagnosis and Treatment. McGraw-Hill; 2025.
\end{enumerate}
\clearpage
